\chapter{Conclusions and further work}
\label{ch:conclusion}

\section{Answers to the research questions}

The training sample contains a negative association after the main controls.
Firm-days with a higher negative-story-share rank tend to have a lower
assigned-window return rank after allowing for mean sentiment and story count.
The baseline coefficient is $-0.00831$ and remains negative after controlling
for recent returns, volatility and recent market sensitivity.

This result provides only exploratory support for H1a. Negative-story share was
chosen from nine candidates using the same training data. The adjusted test
reduces the risk of a chance finding, but confirmation requires a new sample.

FNSPID does not record reliable story-availability times. The association
appears in both the assigned return window and an earlier window. The evidence
therefore shows a relationship within mapped return windows, not one based on
news known before the return began.

The fixed testing-period result does not confirm the training association. During
2020--2023, the unchanged coefficient changes sign to $+0.00583$. The
multi-story estimate is also inconclusive, and no aggregation rule survives
correction. The testing-minus-training contrast persists under longer HAC lags
and a trailing-beta benchmark. Annual estimates are imprecise and do not identify
a break at the chosen boundary. H1b is therefore not supported.

Low power limits what that failure proves. The testing interval still contains
small negative effects, so it cannot establish that the true coefficient is
zero. It does establish that the pre-specified test did not produce the
negative evidence needed to support H1b.

The economic evidence is less ambiguous. The observed-scale effect is small.
The high-minus-low return spread has the sign specified by H2, but its interval
crosses zero. The best break-even cost is $0.625$ basis points per side, versus
the assumed 10 basis points. H3 fails: the training association does not fund
the turnover needed to trade it.

The risk tests provide no economic evidence in support of the trading rule. In
FNSPID, the news rule lowers one downside measure but worsens maximum drawdown.
Because current-session news
sets current-session exposure, this is a same-session check rather than a
strategy formed from information known in advance. Its return advantage over an
after-the-fact constant with the same average exposure is also imprecise.

The LSEG rules fail against controls selected in the training folds with the
same average exposure, and no prompt variant passes its pre-specified criteria. The
samples are too small to show whether the relationship repeats across sources.
H4 therefore fails.

\Cref{tab:hypothesis-verdicts} summarises these conclusions.
\begin{table}[!ht]
\centering
\caption[Hypothesis and transfer verdicts]{Final hypothesis and cross-source verdicts. ``Consistent'' means the selected training-period estimate has the direction stated in H1a, but it is exploratory rather than a confirmation.}
\label{tab:hypothesis-verdicts}
\small
\begin{tabular}{@{}l >{\raggedright\arraybackslash}p{3.3cm} >{\raggedright\arraybackslash}p{8.0cm}@{}}
\toprule
Claim & Verdict & Reason \\
\midrule
H1a & Consistent, not confirmed & Negative selected training-period coefficient; timing is intraday and assignment remains ambiguous. \\
H1b & Not supported & Neither testing-period conditional coefficient is negative, and no testing-period rule survives correction. \\
H2 & Direction only & High-minus-low spread is negative, descriptive and imprecise. \\
H3 & Fails & Turnover makes the daily implementation negative after the declared cost. \\
H4 & Fails & No LSEG risk rule meets all matched-control economic criteria. \\
Cross-source result & Unresolved & Both LSEG FinBERT intervals cross zero and their MDEs are many times the FNSPID coefficient. \\
\bottomrule
\end{tabular}
\end{table}

\FloatBarrier

\section{Contribution}

The academic contribution is a direct comparison of nine firm-day summaries,
followed by a fixed sequence of stronger checks for the selected measure. The
checks add controls, locate the association within the return window, repeat it
in a fixed testing period and test its economic value. Each one answers a
question that the training coefficient cannot answer alone.

Earlier studies link negative financial news with weaker returns
\citep{tetlock2008words,uhl2014,heston2017}. Other work shows that features
beyond an average can matter, including word placement, sentiment dispersion
and news selection \citep{allee2015,isakin2023,uhl2021}. The training result
agrees with this earlier evidence, but supports only a small, sample-specific
association.

The coefficient weakens and becomes inconclusive on multi-story days. FNSPID
also cannot distinguish separate events from repeated coverage, even though
repeated stories can receive a different market response
\citep{tetlock2011stale}. The evidence therefore does not show that a hard
threshold is generally better than averaging.

The practical contribution is a clear screening process. A credible news signal
must meet three conditions. It must be available before the return begins,
persist in later data and remain useful after costs relative to a control with
similar market exposure. This signal fails all three conditions. Its timing is
uncertain, its sign changes in testing and its return does not cover trading
costs.

This result is consistent with wider evidence that return predictors often
weaken outside their original samples \citep{welch2008,mclean2016}. It also shows
why statistical significance is not enough for a high-turnover strategy, where
trading costs can remove apparent value \citep{korajczyk2004,novymarx2016}.

\section{Limitations}

Missing timing information is the main limitation. Date-based FNSPID mapping can
place a story inside a return window that had already opened, and aggregate
return legs cannot reconstruct absent timestamps. The evidence therefore
supports neither causality nor a trading rule that assumes pre-return
availability.

Negative share was chosen from nine measures in the training period, so its
usual interval understates selection uncertainty. The corrected comparison
helps but does not replace a new sample.

The pre-specified analysis improves transparency, but some testing-period
outcomes had already been examined.

The firm sample also limits generalisation. FNSPID includes only firms with news
coverage in every year, which favours surviving, well-covered firms. The sample
was not further restricted to common shares or another security type. It
therefore does not represent the full US equity market.

Story-count results do not identify a mechanism. Coefficients move towards zero
in denser subsets. These estimates are not directly comparable because the ranks
and controls differ across models. The increase in singleton days is also too
small to explain the positive coefficient.

Measurement remains uncertain. FinBERT's classifications were not compared with
human labels in either source, and class shares change modestly through time.
FNSPID also lacks stable
story-family and publisher identifiers, so repeated reports of one event can
raise negative share. Story count controls coverage volume, not novelty. Both
classifier drift and editorial repetition may therefore contribute to the
observed instability.

The economic calculations simplify implementation. They use proportional costs
and estimated adjusted opens, omitting variation in spreads, order size,
borrowing, price impact and cash interest. Most omissions would further reduce
daily long-short returns. The FNSPID downside result is also sample-specific:
overall drawdown worsens, no signal lag is imposed, and activation does not carry
consistently into LSEG.

\section{Further work}

\subsection{Establish point-in-time news availability}

The next study should first establish the order of news and returns. Every story
needs a verified availability timestamp and a pre-specified trading time. A
fixed mapping can then assign pre-open news to the coming session, intraday news
to a later window and post-close news to the next session. Only after fixing this
timing rule should a new signal be specified.

Reuters/LSEG is the practical starting point because it has richer availability
fields. A new FNSPID dataset with reliable timestamps would be required. The
sources should remain separate because their timing, coverage, density and
return definitions differ.

The new data should also link first reports to later updates, using stable
identifiers and update chains where available. Any automated matching of similar
stories should first be tested against hand-checked pairs before returns are
examined.

Comparing results from distinct-event families with results from raw-story
aggregation would then test whether the current result reflects separate events
or repeated coverage.

\subsection{Validate the threshold and sentiment model}

A human-label audit should cover both sources, several years, confidence levels,
predicted classes and story-count groups. Reviewers should label each story for
the named company and report agreement. This would reveal where FinBERT errors
affect negative share.

Future work could compare a second scorer or a flexible probability model under
rules set before either output is linked to returns. One-story and multi-story days
should remain separate because a binary label and a distribution across stories
measure different things.

\subsection{Test on a new sample with adequate power}

The 33-firm LSEG samples cannot show whether a relationship as small as the
FNSPID result repeats in another source. They could reliably detect only effects
7.6 and 13.6 times larger than the FNSPID coefficient. A new sample should
contain several hundred liquid US firms. Before returns are examined, the
coefficient, interval, simple sort and decision rule should be specified.

Interpretation should depend on both interval width and sign. A wide interval
remains inconclusive. A precise estimate near zero would argue against
the relationship repeating at the relevant scale. A small, precise negative
estimate would support a statistical relationship in both sources, but not
economic value.

Any new trading test must also address the roughly sixteen-fold gap between the
best break-even estimate and the assumed cost. A future risk test could retain
the pressure score, 2-basis-point cost, sentiment-free volatility benchmark and
matched-exposure control, but must form exposure only from news verified as
available beforehand. Corrected downside improvement, a return lower bound,
drawdown and pre-specified stability checks should remain mandatory.

Further tuning on the current data cannot solve the central problems. The study
has not established pre-return availability, the fixed relationship changes sign
in testing, and the economic rules fail after costs. A credible new estimate
requires verified timestamps, story-family identities, audited labels and a
genuinely untouched sample.
